% ============================================================
%  BORRADOR PROPUESTA — Subsección: Dataset y metodología
%  Serie macroeconómica FRED (DGS10) como reemplazo/complemento
%  del dataset NASA HTTP en la Actividad 6.
% ============================================================

% ------------------------------------------------------------
\subsection{Dataset y metodología}
% ------------------------------------------------------------

Se analiza la serie diaria \texttt{DGS10} —\textit{10-Year Treasury Constant
Maturity Rate}— disponible en el repositorio FRED (Federal Reserve Economic
Data) del Federal Reserve Bank of St.\ Louis. Esta tasa representa el retorno
 de los bonos del Tesoro de los Estados Unidos a 10 años, uno de los activos
 de referencia más líquidos del mundo y un barómetro clave de las
 expectativas de política monetaria, inflación y crecimiento económico a largo
 plazo.

La serie objeto de estudio son los niveles diarios de la tasa, interpolados
linealmente para cubrir los 717 valores faltantes correspondientes a días
no hábiles (feriados, cierres de mercado). La interpolación preserva la
dinámica de corto plazo sin introducir saltos artificiales. La serie
resultante contiene $16\,801$ observaciones que abarcan el periodo
2 de enero de 1962 al 26 de mayo de 2026, con media
$\mu = 5{,}8139\,\%$, mediana $5{,}4300\,\%$ y desviación estándar
$\sigma = 2{,}9353\,\%$. El coeficiente de variación ($\approx 0{,}50$)
refleja un régimen de bajas tasas post-crisis (mínimo $0{,}52\,\%$ en 2020)
intercalado con ciclos de tightening (máximo $15{,}84\,\%$ en 1981).

Desde la óptica del análisis fractal, la serie \texttt{DGS10} es
particularmente atractiva por tres razones estructurales. En primer lugar,
la política monetaria moderna opera mediante suavización de tasas de
interés: los bancos centrales ajustan el 
federal funds rate de forma
gradual y predecible, generando una fuerte autocorrelación de largo plazo
en la curva de rendimientos. En segundo lugar, los shocks macroeconómicos
(recesiones, crisis financieras, pandemias) dejan huellas persistentes en
la tasa a 10 años, ya que los mercados de bonos descuentan expectativas de
recuperación a horizontes de varios años. En tercer lugar, la curva de
rendimientos actúa como un sistema de retroalimentación multiescala:
interacciones diarias de arbitraje, decisiones trimestrales de política
monetaria y ciclos de endeudamiento soberano de décadas coexisten en la
misma serie temporal. Estas propiedades son análogas, mutatis mutandis, a
las observadas en la variabilidad cardíaca (Actividades~3 y~5) y en el
tráfico transaccional Cloud (sección anterior): retroalimentación temporal,
interacción multiescala y heterogeneidad estructural.

La enorme autocorrelación de primer orden observada ($\rho_1 = 0{,}9998$)
es indicio inequívoco de memoria persistente. A diferencia de las
 diferencias diarias, cuya autocorrelación decae abruptamente
($\rho_1 = 0{,}0809$), los niveles exhiben una dinámica de
\emph{persistencia} compatible con un proceso de dependencia de rango largo.
Por esta razón, el análisis fractal se aplica directamente sobre la serie
en niveles interpolada, sin transformaciones adicionales, garantizando
comparabilidad metodológica con los dominios biológico y tecnológico
tratados previamente.

Se espera un parámetro de Hurst en el rango $\hat{H} \approx 0{,}65$–$0{,}85$.
Este valor se justifica teóricamente: si la tasa a 10 años respondiera a un
proceso sin memoria ($H=0{,}5$), los cambios diarios serían impredecibles y
la serie de niveles no exhibiría la suavidad observable en los datos. La
literatura empírica sobre tasas de interés soberanas reporta consistentemente
$H>0{,}5$ (véase, por ejemplo, los estudios de dependencia de rango largo
en curvas de rendimientos de Reisen y Cribari-Neto). Un valor cercano a
$0{,}7$ indicaría que las decisiones de política monetaria y las
expectativas de inflación generan correlaciones que decaen lentamente,
mientras que valores superiores a $0{,}8$ sugerirían una rigidez extrema
en la curva, típica de regímenes de 
\emph{forward guidance} explícito.

Los estimadores Varianza–Tiempo y Log–Escala, así como el clasificador
mono/multifractal de la Actividad~4, se aplican sin modificación
metodológica. La continuidad en la cadena de procesamiento es esencial:
permite situar la dinámica de la deuda soberana en el mismo espacio de
parámetros que las señales cardíacas y el tráfico web, revelando si la
fractalidad es un epifenómeno de dominios específicos o una propiedad
emergente de sistemas complejos con retroalimentación persistente.
